\subsection{ACRO}
\label{sec:acro}
\paragraph{Actor--Critic Orchestration.}
We implement the objective in Eq.~\ref{eq:orchastration-advantage}
through an Actor--Critic structure. The Critic evaluates the
success prospects of policy continuations, and the Actor uses
these evaluations to select and execute a repositioning motion.
Directly evaluating the orchestration advantage requires searching
for alternative targets and assessing their continuation values.
Doing so at every policy step would spend time even when the
current execution is likely to succeed. We therefore adopt a
sparse Actor--Critic structure: a Q-based assessment identifies
when to activate the Actor, which then searches for a promising
target using V. Once activated, the current action value is fixed,
so maximizing V over candidate targets also maximizes their
estimated orchestration advantage.

\paragraph{Evaluating Continuations.}
Repositioning and retries consume time that could otherwise be
used to complete the task. A useful continuation must therefore
offer a high probability of success within the available time.
We learn Q and V from rollouts of the deployed VLA by modeling
the time to task completion. Q conditions on the observation
history and proposed action, while V uses the observation history
alone. Both predict completion hazard rates, yielding success
probabilities over a forecast horizon $h$:
\begin{equation}
\widehat V(s)=1-e^{-\int_0^h\lambda_V(u\mid s)\,du},
\qquad
\widehat Q(s,a)=1-e^{-\int_0^h\lambda_Q(u\mid s,a)\,du}.
\label{eq:time-to-success-values}
\end{equation}
Here, $\lambda_V$ and $\lambda_Q$ are the predicted completion
rates conditional on the task not yet being completed. We
parameterize them as piecewise-constant rates over time intervals~\citep{kvamme2021survival}.
The forecast horizon $h$ is the interval over which the Critic predicts success; the episode budget counts all physical steps taken by the policy and by retraction.

These predictions must account for both observed successes and
executions that end before success is observed. A successful
rollout provides a completion time; an incomplete rollout indicates
only that completion did not occur during the observed interval.
We incorporate both through the censored time-to-success loss
\begin{equation}
\ell_{\mathrm{surv}}(\lambda;d,\delta)
=\int_0^d\lambda(u)\,du-\delta\log\lambda(d),
\label{eq:survival-loss}
\end{equation}
where $d$ is the duration from a policy decision to completion or
observation cutoff, and $\delta$ indicates whether completion
was observed. We apply this loss to both Q and V.

The value of a proposed action should also agree with the
continuation it produces. For an action interval of duration
$\Delta t$ without termination, the finite-horizon Bellman relation gives~\citep{sutton2018reinforcement}
\begin{equation}
Q_h^\pi(s,a)
=\mathbb E\!\left[V_{h-\Delta t}^\pi(s')\mid s,a\right].
\label{eq:qv-consistency}
\end{equation}
We encourage this consistency by additionally training Q toward
a stop-gradient V prediction at the next observed state, reducing
the forecast horizon by the elapsed action duration. Observed
success in the final action interval supplies a target of one.
Censored endpoints without a subsequent observation are excluded
from this additional bootstrap term.

Before target search, the alternative continuation value is
unavailable. We therefore use Q to predict whether the current
execution will fail within the task horizon. With the Critic
fixed, we train a logistic classifier on recent Q estimates
using rollout-level failure labels. Its prediction determines
when to activate the Actor. This provides a proxy for identifying
opportunities for positive advantage while reserving target
search and physical repositioning for selected moments.

\paragraph{Value-Guided Repositioning.}
The Actor must find a promising target and reach it with enough
time left for task completion. Searching more broadly may reveal
better targets, but increases search time and exposes selection
to value errors at unfamiliar states. We therefore use configurations
recorded along the active execution trajectory as candidates.
These records provide observations and value estimates for target
selection, together with a path for guiding physical movement.

Among the eligible recorded contexts $\mathcal R$, the Actor
selects the target with the highest estimated V. The recorded
value serves as a proxy for the prospects of resuming from its
associated robot configuration. To reduce movement time, Trajectory
Retraction shortcuts the recorded path while keeping each segment
near the visited positions. The target-selection rule and shortcut
condition are
\begin{equation}
\hat s\in\operatorname*{arg\,max}_{s'\in\mathcal R}\widehat V(s'),
\qquad
\max_{p\in[p_i,p_j]}\operatorname{dist}(p,\mathcal P)\le\rho.
\label{eq:value-guided-repositioning}
\end{equation}
Here, $\mathcal P$ is the set of recorded end-effector positions
along the selected trajectory segment, and $\rho$ is the allowed
deviation. After target selection, we greedily choose the farthest
waypoint toward the target whose connecting segment satisfies
the condition, approximated by sampling points along the segment.
The Actor follows the resulting path, and the VLA receives a
fresh observation of the reached state to generate a new
continuation. Subsequent execution updates the active trajectory
for later orchestration decisions. Algorithm~\ref{alg:acro} summarizes the online orchestration loop.
\begin{algorithm}[t]
\caption{Actor--Critic Rollout Orchestration (ACRO)}
\label{alg:acro}
\small
\begingroup
\algrenewcommand{\algorithmiccomment}[1]{\hfill\textcolor{linkblue}{$\triangleright$ #1}}
\begin{algorithmic}
\State Initialize trajectory history $\mathcal B$, Q history $\mathcal Q$, and elapsed steps $n\gets0$
\While{$n<T$ and the episode has not terminated}
    \State Observe $s$ from the reached scene; sample $a\sim\pi(\cdot\mid s)$
    \State $q\gets\widehat Q(s,a)$, $v\gets\widehat V(s)$ \Comment{Evaluate continuations}
    \State Store $(s,v)$ and the robot pose in $\mathcal B$; append $q$ to $\mathcal Q$
    \If{$g(\mathcal Q)>\tau$ and re-entry is available}
        \State $\hat s\gets\operatorname*{arg\,max}_{s'\in\mathcal R}\widehat V(s')$
            \Comment{Select an eligible recorded target}
        \State $\mathcal P\gets$ recorded path from the current pose toward $\hat s$
        \State Shortcut $\mathcal P$ within distance $\rho$ of visited positions \Comment{Trajectory Retraction}
        \State Follow the shortened path; let $\Delta n$ be the physical steps used
    \Else
        \State Execute $a$ until the next query; let $\Delta n$ be the steps used
    \EndIf
    \State Update the active trajectory $\mathcal B$; $n\gets n+\Delta n$
\EndWhile
\end{algorithmic}
\endgroup
\footnotesize
The budget $T$ counts all physical steps; motion stops at termination or budget exhaustion.
Re-entry availability includes a nonempty
eligible set $\mathcal R$ and the intervention conditions in Appendix~\ref{app:trigger}--\ref{app:retraction}.
\end{algorithm}

